%!TEX root = ../thesis.tex
%*******************************************************************************
%****************************** Chapter 3 **************************************
%*******************************************************************************

\chapter{RELATED WORK AND EXISTING SYSTEMS}
\label{chap:related}

\ifpdf
    \graphicspath{{Chapter3/Figs/Raster/}{Chapter3/Figs/PDF/}{Chapter3/Figs/}}
\else
    \graphicspath{{Chapter3/Figs/Vector/}{Chapter3/Figs/}}
\fi

\nomenclature[z-OCR-VQA]{Scene-text VQA}{VQA in which the answer must be read from the text in the image}
\nomenclature[z-Q-Former]{Q-Former}{Querying Transformer: the image--language connector of BLIP-2}

This chapter reviews the five groups of work the thesis builds on: VQA and scene-text
VQA (Section~\ref{sec:rel-vqa}), vision-language models (Section~\ref{sec:rel-vlm}),
parameter-efficient fine-tuning methods (Section~\ref{sec:rel-peft}), Vietnamese VQA
research (Section~\ref{sec:rel-vivqa}) and hallucination measurement
(Section~\ref{sec:rel-halluc}). Section~\ref{sec:rel-summary} summarises the gaps
G1--G4 stated in Chapter~\ref{chap:intro} and shows which gaps the thesis addresses.

%********************************** Section 3.1 ********************************
\section{Visual question answering and scene-text VQA}
\label{sec:rel-vqa}

\subsection{The VQA task}

Antol et al.~\cite{antol2015vqa} framed VQA as an open-ended task: given an image and
a natural-language question, the system answers with a word or a short phrase. Their
VQA dataset collects ten human answers per question and defines the VQA Accuracy
metric: an answer is fully correct if at least three people gave the same answer,
i.e.\
\begin{equation}
  \mathrm{Acc}(a) = \min\left(\frac{\#\{\text{humans who answered } a\}}{3},\ 1\right).
  \label{eq:vqaacc}
\end{equation}
Systems of that period combined a CNN encoding the image with an LSTM encoding the question and
then \emph{classified} the answer over a fixed set of frequent answers. This
classification setting works well for short, recurring answers (colour, count,
yes/no), but cannot produce an answer that never appeared in the answer set during
training.

\subsection{Scene-text VQA}

Many real questions can only be answered by reading the text in the image: a shop
name, a price, a phone number. Singh et al.~\cite{singh2019textvqa} introduced the
TextVQA dataset for this type of question and showed that ordinary VQA models can
hardly answer them; their LoRRA model adds OCR tokens to the model and allows it to
\emph{copy} an OCR token as the answer. Biten et al.~\cite{biten2019stvqa} built the
ST-VQA dataset with the same goal and proposed the ANLS metric (Average Normalized
Levenshtein Similarity), a character-level similarity between the answer and the gold
answer that does not heavily penalise small OCR reading errors. M4C~\cite{hu2020m4c}
extends the copying idea with a multimodal transformer that generates the answer step
by step, choosing at each step between a vocabulary word and an OCR token.

The common lesson of this line of work is that, in scene-text VQA, the quality of the
answer is bounded by the quality of the text that is read. If OCR misreads the text or
sorts the text regions in the wrong order, the model cannot copy the correct answer.
This is why the thesis measures the quality of Vietnamese OCR and the reading order
separately (Chapter~\ref{chap:methods}).

%********************************** Section 3.2 ********************************
\section{Vision-language models}
\label{sec:rel-vlm}

Current vision-language models (VLMs) share a common template: an image encoder turns
the image into a sequence of vectors, a connector maps these vectors into the space of
the language model, and a large language model (LLM) reads the image tokens together
with the question to generate the answer. VLMs differ in the connector, in how they
handle image resolution and in their training data. Table~\ref{tab:rel-vlm} compares
the four VLMs that the project proposal allows as baselines.

\paragraph{BLIP-2.} Li et al.~\cite{li2023blip2} keep both the image encoder and the
LLM frozen and train only a small module in between, the Q-Former. The Q-Former uses a
small number of learned queries to extract information from the image into a fixed
sequence of tokens fed to the LLM. This is cheap to train, but the small, fixed number
of image tokens makes it hard for the model to keep small details such as text in the
image.

\paragraph{LLaVA and LLaVA-1.5.} Liu et al.~\cite{liu2023llava} connect the image
encoder to the LLM with a single projection layer and train on conversational data
about images (visual instruction tuning). LLaVA-1.5~\cite{liu2024llava15} replaces the
projection with an MLP, increases the image resolution and adds a prompt asking for
short answers; this prompt alone is enough for the model to answer in the short style
that VQA datasets require. LLaVA-1.5 also releases both a fully fine-tuned version and
a LoRA version (Section~\ref{sec:rel-peft}).

\paragraph{Florence-2.} Xiao et al.~\cite{xiao2024florence2} build a unified
encoder--decoder vision model that takes text prompts and produces results for
captioning, object detection, grounding and OCR. The model is small (0.23 and
0.77 billion parameters), so it is easy to train, but it is not a multilingual
conversational model and was not trained for open-ended question answering in
Vietnamese.

\paragraph{Qwen2.5-VL.} Bai et al.~\cite{bai2025qwen25vl} use a ViT encoder that takes
images at their native resolution (native dynamic resolution): the image is cut into
14×14-pixel patches, and the connector merges each group of 2×2 patches into one token,
so the number of image tokens is proportional to the image size. The language model is
the multilingual Qwen2.5, and the model is trained extensively on documents and text in
images; the 3-billion-parameter version reaches 79.3 on TextVQA according to the
authors. The thesis chooses Qwen2.5-VL-3B-Instruct as the baseline; the detailed
architecture and the reasons for the choice are presented in Chapter~\ref{chap:methods}.

\begin{table}[htbp]
\caption{Comparison of the four VLMs within the scope of the proposal}
\label{tab:rel-vlm}
\centering
\small
\begin{tabular}{p{0.13\textwidth}p{0.25\textwidth}p{0.24\textwidth}p{0.26\textwidth}}
\toprule
\textbf{Model} & \textbf{Image--language connector} & \textbf{Strengths for this project} & \textbf{Weaknesses for this project} \\
\midrule
BLIP-2 & Q-Former, fixed number of image tokens; image encoder and LLM frozen & Cheap to train; basis of ViTextBLIP-2 & Few image tokens, easily loses small text details \\
LLaVA-1.5 & MLP projecting each patch into the LLM & Simple architecture; evidence that LoRA matches full fine-tuning & Mainly English; fixed image resolution \\
Florence-2 & Unified encoder--decoder with text prompts & Very small; has a dedicated OCR task & Not a multilingual conversational model \\
Qwen2.5-VL & Native-resolution ViT, 2×2 patches merged into one token & Multilingual; reads text in images well; has a 3B version & Number of image tokens grows with image size, memory-hungry \\
\bottomrule
\end{tabular}
\end{table}

%********************************** Section 3.3 ********************************
\section{Parameter-efficient fine-tuning}
\label{sec:rel-peft}

\subsection{Why PEFT is needed}

Fully fine-tuning a large model for each task means updating and storing a copy of
every parameter. For the 175-billion-parameter GPT-3, each copy takes about 350~GB in
16 bits, and training also has to keep gradients and optimiser
states~\cite{hu2022lora}. Parameter-efficient fine-tuning (PEFT) keeps the original
weights frozen and trains only a small number of added parameters. Early
approaches include adapters~\cite{houlsby2019adapter}, which insert small layers between
the layers of the network, and prefix-tuning~\cite{li2021prefix}, which learns a
sequence of vectors placed before the input. Adapters add inference latency because
they add sequential layers; prefix-tuning takes up part of the sequence length and is
hard to optimise.

\subsection{LoRA}

Hu et al.~\cite{hu2022lora} assume that the change in the weights during adaptation
has a low rank. For a weight matrix $W_0 \in \mathbb{R}^{d \times k}$, LoRA keeps $W_0$
fixed and learns two small matrices $B \in \mathbb{R}^{d \times r}$ and
$A \in \mathbb{R}^{r \times k}$ with $r \ll \min(d, k)$:
\begin{equation}
  h = W_0 x + \frac{\alpha}{r} B A x .
  \label{eq:lora}
\end{equation}
$B$ is initialised to zero, so at the start the model keeps its original behaviour.
After training, $BA$ can be added directly to $W_0$, so inference is not slower. On
GPT-3, LoRA matches or exceeds full fine-tuning with about 10,000 times fewer trainable
parameters and about 3 times less GPU memory.

\subsection{QLoRA}

LoRA reduces the number of parameters to learn but still has to keep the original
weights in 16 bits in memory. Dettmers et al.~\cite{dettmers2023qlora} quantise the
original weights to 4 bits and train LoRA on top, with three techniques: the 4-bit
NormalFloat (NF4) data type, designed for weights with a near-normal distribution;
double quantisation, which also compresses the quantisation constants; and paged
optimisers, which avoid running out of memory during consumption peaks. As a result,
a 65-billion-parameter LLM can be fine-tuned on one 48~GB GPU, while the results match
16-bit fine-tuning in the authors' experiments. The paper also recommends attaching LoRA
to every linear layer to compensate for the quantisation error. The thesis uses QLoRA
because it is the only way to fine-tune a VLM with several billion parameters on one
Colab GPU; the choice between attaching LoRA only to the attention matrices or to every
linear layer is tested with an ablation in Chapter~\ref{chap:results}.

\subsection{PEFT for vision-language models}

Sung et al.~\cite{sung2022vladapter} compare PEFT methods in a controlled way on
image--language tasks, including VQA, and show that a shared adapter with only 4.18\%
trainable parameters comes close to full fine-tuning. LLaVA-1.5 gives similar evidence
at a larger scale: with the 13-billion-parameter model, the LoRA version reaches 80.0 on
VQAv2 (the same as full fine-tuning), 60.2 against 61.3 on TextVQA, and 86.7 against
85.9 F1 on POPE~\cite{liu2024llava15}. Table~\ref{tab:rel-peft} summarises the
adaptation methods.

\begin{table}[htbp]
\caption{Ways to adapt a large model}
\label{tab:rel-peft}
\centering
\small
\begin{tabular}{p{0.20\textwidth}p{0.22\textwidth}p{0.20\textwidth}p{0.26\textwidth}}
\toprule
\textbf{Method} & \textbf{Trainable parameters} & \textbf{Inference latency} & \textbf{Training memory} \\
\midrule
Full fine-tuning & All & Unchanged & Very large: weights, gradients, optimiser states \\
Adapter & A few percent & Higher (extra layers) & Keeps 16-bit weights \\
Prefix-tuning & Very few & Slightly higher (longer sequence) & Keeps 16-bit weights \\
LoRA & Below 1\% & Unchanged after merging & Keeps 16-bit weights \\
QLoRA & Below 1\% & Unchanged after merging & 4-bit weights; the smallest \\
\bottomrule
\end{tabular}
\end{table}

However, all of this evidence comes from English. Among the works we reviewed,
none verifies LoRA or QLoRA on a VLM for Vietnamese VQA or for Vietnamese scene-text
VQA. This is gap G3, on which the thesis focuses.

%********************************** Section 3.4 ********************************
\section{Vietnamese VQA}
\label{sec:rel-vivqa}

\subsection{Datasets}

\paragraph{ViVQA.} Tran et al.~\cite{tran2021vivqa} released the first Vietnamese VQA
dataset, with 10,328 images and 15,000 question--answer pairs. The answers are short,
mostly one to three words, so the task is usually framed as classification.

\paragraph{OpenViVQA.} Nguyen et al.~\cite{nguyen2023openvivqa} built a dataset with
open-ended, longer answers that must be generated rather than chosen from a fixed set;
about half of the samples involve text in the image~\cite{nguyen2024vitextvqa}.

\paragraph{ViTextVQA.} Nguyen et al.~\cite{nguyen2024vitextvqa} released the first
large-scale Vietnamese scene-text VQA dataset, with 16,762 images and 50,342 questions,
split into train, dev and test in a 7:1:2 ratio. The answers are usually text copied
from the image, such as shop names, slogans or prices. Humans reach 88.00 EM and 96.42
F1 on this dataset. The authors also show that sorting the OCR text from left to right
and from top to bottom helps the models answer better. The thesis uses ViVQA and
ViTextVQA; detailed statistics are presented in Chapter~\ref{chap:methods}.

\subsection{Methods}

Table~\ref{tab:rel-vivqa} summarises representative methods over time.

\begin{table}[htbp]
\caption{Representative Vietnamese VQA studies}
\label{tab:rel-vivqa}
\centering
\small
\begin{tabular}{p{0.21\textwidth}p{0.33\textwidth}p{0.16\textwidth}p{0.20\textwidth}}
\toprule
\textbf{Study} & \textbf{Approach} & \textbf{Published result} & \textbf{Limitation} \\
\midrule
ViVQA~\cite{tran2021vivqa}, 2021 & LSTM, Bi-LSTM with several word embeddings; Hierarchical Co-Attention; answer chosen from a fixed set & Acc 34.96 on ViVQA & Hand-crafted features; cannot read text \\
Bi-directional Cross-Attention~\cite{nguyentran2022bica}, 2022 & Pre-trained encoders, bidirectional cross-attention between image and question & Beats the ViVQA baseline & Still a classification task \\
PAT~\cite{nguyen2023pat}, 2023 & Parallel Attention Transformer, hierarchical question features & Acc 60.55 on ViVQA & Classification; does not read text \\
BARTPhoBEiT~\cite{tran2023bartphobeit}, 2023 & BARTpho for Vietnamese combined with a BEiT-style image encoder & Acc 68.58 on ViVQA & Classification; does not read text \\
AViVQA-TranConI~\cite{nguyen2024tranconi}, 2024 & BLIP-2 (global features) and EfficientNet (local features), BARTpho, fused with BEiT-3; classification over 353 answers & Acc 71.04 on ViVQA & Cannot generate new answers; no OCR \\
ViTextBLIP-2~\cite{nguyen2024vitextvqa}, 2026 & SwinTextSpotter reads the text, a Q-Former connects image and OCR text, a frozen ViT5 generates the answer & EM 25.48; F1 53.95 on ViTextVQA & Depends on external OCR; far from humans (88.00 EM) \\
\bottomrule
\end{tabular}
\end{table}

These studies follow a clear path: from hand-crafted features with attention (2021), to
pre-trained encoders and multimodal fusion modules (2022--2024), to a generative model
based on BLIP-2 for data with text (2025). Two problems recur. First, most works on
ViVQA treat VQA as classification over a fixed answer set, so they cannot produce
answers unseen during training, which is mandatory for questions that require reading.
Second, reading Vietnamese text with diacritics in images remains a weak point:
ViTextBLIP-2 relies on an external OCR system kept fixed during training, so the OCR's
reading errors pass straight into the answer.

A caution about comparisons: the Accuracy of the classification methods on ViVQA is
computed over a fixed answer set, which is not the same scoring as the EM of the
free-generation model used in this thesis. The ViVQA numbers in
Table~\ref{tab:rel-vivqa} are therefore for reference only. On ViTextVQA, the thesis
compares directly with ViTextBLIP-2 because both use the same test set and the same EM
and F1 metrics (Chapter~\ref{chap:results}).

%********************************** Section 3.5 ********************************
\section{Hallucination and how to measure it}
\label{sec:rel-halluc}

Hallucination in VLMs is the phenomenon in which the model generates fluent content
that is not in the image, for example mentioning an object that does not appear.
Rohrbach et al.~\cite{rohrbach2018chair} measure it on image captioning with the CHAIR
metric, which counts the share of mentioned objects that are not in the image labels.
This measure depends on extracting object names from the sentence and on the sentence
length.

POPE~\cite{li2023pope} switches to yes/no questions: for each image it asks ``Is there
\emph{an X} in the image?'' for both objects that are present and objects that are not.
The absent objects are chosen in three ways: at random, among the objects most frequent
in the whole dataset, and among the objects that often co-occur with the objects in the
image (adversarial). Results are reported as Accuracy, Precision, Recall, F1 and the
``yes'' ratio. Because the answer is only yes or no, POPE needs no sentence analysis and
is easy to apply to other languages. However, no POPE-style benchmark exists in
Vietnamese (gap G4). The thesis builds POPE-vi following exactly this protocol to test
hypothesis H4 (Chapter~\ref{chap:discussion}).

%********************************** Section 3.6 ********************************
\section{Summary and research gaps}
\label{sec:rel-summary}

Table~\ref{tab:rel-gaps} links each gap stated in Chapter~\ref{chap:intro} to the
evidence from the works above and to how the thesis addresses it.

\begin{table}[htbp]
\caption{Gaps, evidence and how the thesis addresses them}
\label{tab:rel-gaps}
\centering
\small
\begin{tabular}{p{0.20\textwidth}p{0.36\textwidth}p{0.34\textwidth}}
\toprule
\textbf{Gap} & \textbf{Evidence} & \textbf{How the thesis addresses it} \\
\midrule
G1. English-centric pre-training & VLMs are trained mainly on English; LLaVA-1.5 shows only emergent multilingual ability, not dedicated training & Choose a multilingual baseline; measure wrong-language and wrong-format errors in the zero-shot setting \\
G2. Weak Vietnamese VQA, especially for reading text & ViTextBLIP-2 reaches 25.48 EM against 88.00 for humans; the ViVQA methods are classifiers that do not read text & Free-form answer generation; Vietnamese OCR text with diacritics, in reading order, in the prompt \\
G3. No PEFT for Vietnamese VLMs & LoRA, QLoRA, VL-Adapter and LLaVA-1.5 are only verified on English & QLoRA on Qwen2.5-VL-3B; OCR-aware QLoRA as the main contribution \\
G4. No Vietnamese hallucination benchmark & POPE and CHAIR exist only in English & Build POPE-vi; check whether QLoRA increases hallucination \\
\bottomrule
\end{tabular}
\end{table}

In short, previous work shows that VLMs can read text in images, that PEFT is enough to
adapt VLMs, and that the order of the OCR text matters for Vietnamese VQA. What is
missing is a study that combines these three ideas for Vietnamese under a small resource
budget, with a reproduced baseline, ablations and hallucination measurement.
Chapter~\ref{chap:methods} presents how the thesis does this.
